% *********************************************************************
% © 2021 Jeremy Sylvestre
%
% Permission is granted to copy, distribute and/or modify this document
% under the terms of the GNU Free Documentation License, Version 1.3 or
% any later version published by the Free Software Foundation; with no
% Invariant Sections, no Front-Cover Texts, and no Back-Cover Texts. A
% copy of the license is included in the appendix entitled “GNU Free
% Documentation License” that appears in the output document of this
% PreTeXt source code. All trademarks™ are the registered® marks of
% their respective owners.
%
% *********************************************************************

\newcommand*{\xmin}{-0.5}
\newcommand*{\xmax}{6.5}
\newcommand*{\ymin}{-0.5}
\newcommand*{\ymax}{6.5}

%
% to maintain all segments at same length L
% (in this case L = 0.2)
% use
%   delta-x = (L/2) / sqrt(1 + m^2)
%   delta-y = m * (L/2) / sqrt(1 + m^2)
% (the L/2 is because we draw two segments, each from centre out to either right or left)
%

\begin{tikzpicture}

	% axes
	\draw[->,thick] (\xmin,0) -- (\xmax,0) node[below] {$t$};
	\draw[->,thick] (0,\ymin) -- (0,\ymax) node[right] {$q(t)$};

	\foreach \x in {0,0.5,1,...,6} {
		\foreach \y in {0,0.5,1,...,6} {
			\draw (\x,\y) -- ++($({0.1/sqrt(1+\y*\y/4)},{-0.1*\y/sqrt(1+\y*\y/4)/2})$);
			\draw (\x,\y) -- ++($({-0.1/sqrt(1+\y*\y/4)},{0.1*\y/sqrt(1+\y*\y/4)/2})$);
		}
	}

\end{tikzpicture}
